\documentclass[12pt]{article}
\usepackage[utf8]{inputenc}
\usepackage{bbm}
\usepackage{geometry}
\geometry{a4paper,scale=0.75}
\linespread{1.5}
\usepackage{graphicx} 
\usepackage{float} 
\usepackage{subcaption} 
\usepackage{enumerate}
\usepackage{amsmath}
\usepackage{array}
\newcolumntype{P}[1]{>{\raggedright\arraybackslash}p{#1}}
\usepackage{booktabs}
\usepackage{multirow}
\usepackage{amsfonts}
\usepackage[english]{babel}
\usepackage{amsthm}
\usepackage{dcolumn}
\usepackage{multicol}
\usepackage{stfloats}
\usepackage{placeins}
\usepackage{lscape}
\usepackage[figuresright]{rotating}
\RequirePackage{pdflscape}
\usepackage[toc,page]{appendix}
\usepackage{geometry}
\usepackage{longtable}
\usepackage{comment}
\usepackage{xcolor}
% \usepackage{amsmath}
\usepackage{amssymb}
\usepackage{empheq}
\usepackage{newtxtext,newtxmath}

% -------- enumerated sub-labels (a), (b), … --
\usepackage{enumitem}
\setlist[enumerate,1]{label=(\alph*),ref=\alph*}
% ---------------------------------------------

\usepackage{hyperref}
\hypersetup{hidelinks,
	colorlinks=true,
	allcolors=black,
	pdfstartview=Fit,
	breaklinks=true}
\usepackage{csquotes}
\usepackage{natbib}
\newtheorem{definition}{Definition}
\newtheorem{theorem}{Theorem}
%\newtheorem{proposition}[theorem]{Proposition}
\newtheorem{proposition}{Proposition}[section]
%\newtheorem{lemma}[theorem]{Lemma}
\newtheorem{lemma}{Lemma}[section]
%\newtheorem{corollary}[theorem]{Corollary}
\newtheorem{corollary}{Corollary}[section]
\newtheorem*{remark}{Remark}
\newtheorem{example}{Example}
\newtheorem{exercise}{Exercise}[section]
\numberwithin{equation}{section}
\newtheorem{assumption}{Assumption}[section] % number within sections

\title{Notes on \cite{christianoNominalRigiditiesDynamic2005}}
\author{Harlly Zhou \\ \texttt{hzhouci@ust.hk}}
\date{\today}

\begin{document}
\maketitle

\paragraph{Motivation} Observed inflation inertia and aggregate quantity persistence.

\paragraph{Research qeustion} Can models with moderate degrees of nominal rigidities generate inertial inflation and persistent output movements in response to a monetary policy shock?

\paragraph{Consequences of MP shock} After an expansionary monetary policy shock:
\begin{enumerate}
    \item output, consumption, and investment respond in a hump-shaped fashuon, peaking after about one and a half years and returning to preshock levels after about three years;
    \item inflation responds in a hump-shaped fashion, peaking after about two years;
    \item the interest rate falls for roughly one year;
    \item real profits, real wages and labour productivity rise; and
    \item the growth rate of money rises immediately.
\end{enumerate}

\paragraph{Model} 
\begin{enumerate}
    \item Firm, still, is doing recursive optimization problem with continuation value: Choose its price to maximize the expected continuation value under this price.
    \item Money supply is the sum of cash balance and debt.
\end{enumerate}



\bibliographystyle{../draft/aer}
\bibliography{../draft/network_uncertainty}

\end{document}
